\documentclass[11pt]{article}

\usepackage[margin=1in]{geometry}
\usepackage{amsmath,amssymb,amsthm}
\usepackage{booktabs}
\usepackage{hyperref}
\usepackage{algorithm}
\usepackage{algpseudocode}

\newtheorem{definition}{Definition}
\newtheorem{proposition}{Proposition}
\newtheorem{remark}{Remark}
\newtheorem{corollary}{Corollary}

\title{Inverse Design of Cylindrical Compression Springs\\
from a Target Rate and Fatigue Safety Factor}
\author{Derived from the \texttt{CompressionSpringInverseDesigner} implementation}
\date{}

\begin{document}
\maketitle

\begin{abstract}
We formalize an inverse-design procedure for standard cylindrical helical
compression springs. Given a linear rate target — specified either by two
exact length/force points or by a least-squares fit to several noisy
length/force samples — together with a target fatigue safety factor, the
procedure recovers a wire diameter, mean coil diameter, number of coils,
and free length that reproduce the rate exactly while matching the safety
factor as closely as possible. The search exploits two structural facts:
(i) for a fixed pair of load points the wire stress depends only on the
wire and mean diameters, so the number of coils can always be chosen
after the fact to hit the rate exactly; and (ii) for fixed wire diameter,
the fatigue safety factor is monotonic in the mean diameter away from the
spring-index singularity at $C=4$, so it may be inverted by one-dimensional
root finding. The outer loop is a brute-force scan over a discrete
standard wire-diameter series, each of whose candidates is validated
against buildability constraints (minimum coil count, non-bottoming solid
length, non-interfering pitch) before selection.
\end{abstract}

\tableofcontents

\section{Problem statement}

\subsection{Design requirements}

The design brief consists of a material, a target fatigue safety factor
$n_f^{\star}$, and a target linear rate specified in one of two
equivalent forms:

\begin{itemize}
\item \textbf{Two-point form:} two distinct (length, force) pairs
$(\ell_1, F_1)$, $(\ell_2, F_2)$ on the spring's pre-contact linear
regime;
\item \textbf{Sample form:} a set of $m\ge 2$ measured pairs
$\{(\ell_i, F_i)\}_{i=1}^m$, assumed noisy but linear, from which the
rate is estimated by least squares.
\end{itemize}

\subsection{Design variables and target quantities}

The free geometric variables are the wire diameter $d$, the mean coil
diameter $D$, the number of active coils $n_a$ (equivalently, total
coils $n = n_a + n_0$ for an end/forming--dependent offset $n_0$), and
the free length $L_0$. The design must reproduce a target spring rate
$k^{\star}$ and free length $L_0^{\star}$ exactly, while bringing the
fatigue safety factor $n_f(d,D)$, evaluated at the two extreme load
points, as close as possible to $n_f^{\star}$.

\section{Recovering the rate target}

\subsection{Exact two-point solve}

A compression spring's pre-contact force--deflection relation is linear,
\begin{equation}
F(\ell) = k\,(L_0 - \ell),
\label{eq:linear-law}
\end{equation}
where $\ell$ is the compressed length and $L_0$ the free length. Given
two distinct points $(\ell_1,F_1)$, $(\ell_2,F_2)$ with, without loss of
generality, $\ell_1 < \ell_2$ (so $F_1 > F_2$ for a valid compression
spring), Eq.~\eqref{eq:linear-law} is fully determined:
\begin{equation}
k = \frac{F_1 - F_2}{\ell_2 - \ell_1},
\label{eq:rate-two-point}
\end{equation}
\begin{equation}
L_0 = \ell_1 + \frac{F_1}{k}.
\label{eq:free-length-two-point}
\end{equation}
Equation~\eqref{eq:free-length-two-point} is the zero-force extrapolation
of the line back from the more-compressed (higher-force) point. We write
$\ell_{\mathrm{lo}} = \ell_1$, $F_{\mathrm{hi}}=F_1$,
$\ell_{\mathrm{hi}}=\ell_2$, $F_{\mathrm{lo}}=F_2$ for the ordered pair
used later as the two representative operating points.

\subsection{Least-squares solve from multiple samples}

Given $m \ge 2$ samples $\{(\ell_i,F_i)\}$ with at least two distinct
lengths, fit
\begin{equation}
F(\ell) \approx a\,\ell + b
\label{eq:ls-line}
\end{equation}
by ordinary least squares,
\begin{equation}
(a,b) = \operatorname*{arg\,min}_{a,b}\ \sum_{i=1}^{m}\big(F_i - a\ell_i - b\big)^2 .
\end{equation}
A valid compression-spring fit requires $a<0$; the rate and free length
are then read off the fitted line exactly as in the two-point case,
\begin{equation}
k = -a, \qquad L_0 = -\frac{b}{a}.
\label{eq:rate-ls}
\end{equation}
The two representative operating points are taken at the extreme sampled
lengths $\ell_{\mathrm{lo}}=\min_i \ell_i$, $\ell_{\mathrm{hi}}=\max_i \ell_i$,
using the \emph{fitted} forces
\begin{equation}
F_{\mathrm{hi}} = a\,\ell_{\mathrm{lo}} + b, \qquad
F_{\mathrm{lo}} = a\,\ell_{\mathrm{hi}} + b,
\label{eq:ls-endpoints}
\end{equation}
rather than the raw measured values, so downstream stress checks stay
consistent with the fitted line.

\begin{remark}
For $m=2$, Eq.~\eqref{eq:ls-line} passes through both points exactly, so
Eqs.~\eqref{eq:rate-ls}--\eqref{eq:ls-endpoints} reduce to
Eqs.~\eqref{eq:rate-two-point}--\eqref{eq:free-length-two-point}: the
least-squares path strictly generalizes the two-point path.
\end{remark}

\section{Stress model}

\subsection{Wahl factor and shear stress}

For spring index $C = D/d$, the Wahl correction factor is
\begin{equation}
K_W(C) = \frac{4C-1}{4C-4} + \frac{0.615}{C},
\label{eq:wahl}
\end{equation}
and the corrected shear stress in a round wire under axial load $F$ is
\begin{equation}
\tau(D,d,F) = \frac{8\,D\,F}{\pi d^{3}}\,K_W\!\left(\frac{D}{d}\right).
\label{eq:stress}
\end{equation}
Equation~\eqref{eq:stress} depends on the geometry only through $(D,d)$
and is evaluated at the two representative operating loads
$F_{\mathrm{hi}}, F_{\mathrm{lo}}$ to obtain the operating stress range
\begin{equation}
\tau_{\mathrm{hi}} = \tau(D,d,F_{\mathrm{hi}}), \qquad
\tau_{\mathrm{lo}} = \tau(D,d,F_{\mathrm{lo}}).
\label{eq:stress-extremes}
\end{equation}

\subsection{Fatigue safety factor}

Given a Goodman fatigue model for the material (endurance limit,
ultimate strength, cycle count, and optional shot-peening correction),
the safety factor is a function
\begin{equation}
n_f = n_f(\tau_{\mathrm{hi}}, \tau_{\mathrm{lo}}) ,
\label{eq:safety-factor-def}
\end{equation}
computed from the mean and alternating components
$\tau_m = (\tau_{\mathrm{hi}}+\tau_{\mathrm{lo}})/2$,
$\tau_a = (\tau_{\mathrm{hi}}-\tau_{\mathrm{lo}})/2$
against the modified-Goodman line, as in the companion compression-spring
model. Composing with Eq.~\eqref{eq:stress-extremes} gives the safety
factor purely as a function of the two free diameters,
\begin{equation}
n_f = n_f(D,d) .
\label{eq:safety-factor-Dd}
\end{equation}

\section{Coil count from the rate target}

\subsection{Closed-form active-coil count}

The classical helical-spring rate formula,
\begin{equation}
k = \frac{G\,d^{4}}{8\,D^{3}\,n_a},
\label{eq:rate-formula}
\end{equation}
with $G$ the material shear modulus, is solved for the number of active
coils $n_a$ once $(D,d)$ and the target rate $k^{\star}$
(Eqs.~\eqref{eq:rate-two-point} or~\eqref{eq:rate-ls}) are fixed:
\begin{equation}
n_a(D,d) = \frac{G\,d^{4}}{8\,D^{3}\,k^{\star}}.
\label{eq:active-coils}
\end{equation}

\begin{proposition}
For any admissible pair $(D,d)$ with $D,d>0$, Eq.~\eqref{eq:active-coils}
supplies a coil count that reproduces the target rate exactly. Hence the
rate constraint never restricts the search over $(D,d)$: it is satisfied
identically by the choice of $n_a$.
\end{proposition}

\subsection{Total coil count}

The total coil count adds an end/forming-dependent constant offset
$n_0$ (encoding, e.g., closed-and-ground versus plain ends):
\begin{equation}
n = n_a(D,d) + n_0 .
\label{eq:total-coils}
\end{equation}

\section{Reduction to a one-dimensional search}
\label{sec:reduction}

Because $n_a$ is fully determined by $(D,d)$ through
Eq.~\eqref{eq:active-coils}, the design problem reduces to choosing
$(D,d)$ so that
\begin{equation}
n_f(D,d) = n_f^{\star}
\label{eq:target-equation}
\end{equation}
as closely as possible, subject to buildability constraints
(Section~\ref{sec:constraints}). The wire diameter $d$ is restricted to a
discrete standard series $\mathcal{D}_{\mathrm{std}}$ (e.g.\ from a
diameter/tolerance table), so the two-dimensional search over $(D,d)$
decomposes into
\[
|\mathcal{D}_{\mathrm{std}} \cap [d_{\min}, d_{\max}]|
\]
independent one-dimensional root-finding problems in $D$, one per
candidate $d$.

\begin{proposition}[Monotonicity]
\label{prop:monotonic}
For fixed $d$, the stress $\tau(D,d,F)$ of Eq.~\eqref{eq:stress} is
increasing in $D$ away from the singular spring index $C=4$ (where $K_W$
diverges), since both the explicit factor $D$ and the Wahl factor
$K_W(D/d)$ increase with $D$ in this regime. Consequently the safety
factor $n_f(D,d)$, which decreases with increasing operating stresses, is
monotonically decreasing in $D$ over any spring-index interval bounded
away from $C=4$.
\end{proposition}

\begin{corollary}
On an index bracket $[C_{\min},C_{\max}]$ not containing $C=4$, with
corresponding diameter bracket $[D_{\min},D_{\max}] = [C_{\min}d,\,
C_{\max}d]$, Eq.~\eqref{eq:target-equation} may be solved for $D$ by
standard bracketed root-finding (e.g.\ Brent's method) whenever
\begin{equation}
\big(n_f(D_{\min},d) - n_f^{\star}\big)\big(n_f(D_{\max},d) - n_f^{\star}\big) \le 0 .
\label{eq:bracket-condition}
\end{equation}
\end{corollary}

\begin{remark}
When condition~\eqref{eq:bracket-condition} fails, the target safety
factor is unreachable for this $d$ within the allowed index range; the
nearer of the two bracket endpoints is reported instead, i.e.
\begin{equation}
D^{*}(d) = \operatorname*{arg\,min}_{D\in\{D_{\min},D_{\max}\}}
\big|n_f(D,d) - n_f^{\star}\big| .
\end{equation}
\end{remark}

\section{Buildability constraints}
\label{sec:constraints}

A candidate $(D,d)$, with implied $n_a$ from Eq.~\eqref{eq:active-coils}
and $n$ from Eq.~\eqref{eq:total-coils}, is rejected unless all of the
following hold.

\paragraph{Minimum coil count.}
\begin{equation}
n_a \ge n_a^{\min}.
\label{eq:min-coils}
\end{equation}

\paragraph{Non-bottoming solid length.}
The fully closed (solid) length,
\begin{equation}
L_s = n\,d,
\label{eq:solid-length-inv}
\end{equation}
must not reach or exceed the shortest working length in the target
duty cycle:
\begin{equation}
L_s < \ell_{\mathrm{lo}}.
\label{eq:no-bottom}
\end{equation}

\paragraph{Non-interfering pitch.}
The nominal coil pitch at free length,
\begin{equation}
p = \frac{L_0^{\star}}{n},
\label{eq:pitch}
\end{equation}
must exceed the wire diameter, so that coils do not already touch at
free length:
\begin{equation}
p > d.
\label{eq:pitch-constraint}
\end{equation}

A candidate satisfying Eqs.~\eqref{eq:min-coils},
\eqref{eq:no-bottom}, and~\eqref{eq:pitch-constraint} is \emph{valid}.

\section{Selection among candidates}

\subsection{Per-diameter candidate}

For each standard wire diameter $d\in\mathcal{D}_{\mathrm{std}}$ within
the allowed bounds, Section~\ref{sec:reduction} yields a candidate mean
diameter $D^{*}(d)$, spring index $C^{*}(d)=D^{*}(d)/d$, and achieved
safety factor $n_f^{*}(d) = n_f(D^{*}(d),d)$, together with a validity
flag from Section~\ref{sec:constraints}.

\subsection{Selection rule}

Let $\mathcal{V}\subseteq\mathcal{D}_{\mathrm{std}}$ be the subset of
diameters yielding a valid candidate, and let
\begin{equation}
C_{\mathrm{mid}} = \frac{C_{\min}+C_{\max}}{2}
\end{equation}
be the midpoint of the allowed spring-index range. The selected design
minimizes, lexicographically,
\begin{equation}
d^{\star} = \operatorname*{arg\,min}_{d \in \mathcal{V}}
\Big(\ \big|n_f^{*}(d) - n_f^{\star}\big|,\ \ \big|C^{*}(d) - C_{\mathrm{mid}}\big|\ \Big),
\label{eq:selection-rule}
\end{equation}
i.e.\ closeness to the target safety factor is the primary criterion,
and closeness of the spring index to the middle of its admissible range
breaks (near-)ties.

\begin{remark}
If $\mathcal{V}=\varnothing$ — no standard wire diameter admits a valid
geometry within the allowed spring-index and wire-diameter bounds for
this rate/safety-factor target — the design problem has no feasible
solution under the stated bounds, and the bounds must be relaxed.
\end{remark}

\section{Algorithm summary}

\begin{algorithm}[h]
\caption{Inverse compression-spring design}
\begin{algorithmic}[1]
\Require material, $n_f^{\star}$, rate target data, bounds $[C_{\min},C_{\max}]$, $[d_{\min},d_{\max}]$
\State Compute $(\ell_{\mathrm{lo}}, F_{\mathrm{hi}}, \ell_{\mathrm{hi}}, F_{\mathrm{lo}}, k^{\star}, L_0^{\star})$ via Eqs.~\eqref{eq:rate-two-point}--\eqref{eq:free-length-two-point} or~\eqref{eq:rate-ls}--\eqref{eq:ls-endpoints}
\For{each $d \in \mathcal{D}_{\mathrm{std}} \cap [d_{\min},d_{\max}]$}
    \State Set bracket $[D_{\min},D_{\max}] = [C_{\min}d, C_{\max}d]$
    \State Solve Eq.~\eqref{eq:target-equation} for $D^{*}(d)$ by root finding (Eq.~\eqref{eq:bracket-condition}) or take the nearer bound
    \State Compute $n_f^{*}(d)$, $C^{*}(d) = D^{*}(d)/d$
    \State Compute $n_a$ (Eq.~\eqref{eq:active-coils}), $n$ (Eq.~\eqref{eq:total-coils})
    \State Validate against Eqs.~\eqref{eq:min-coils},~\eqref{eq:no-bottom},~\eqref{eq:pitch-constraint}
\EndFor
\State Select $d^{\star}$ by Eq.~\eqref{eq:selection-rule} among valid candidates
\State \Return geometry $(d^{\star}, D^{*}(d^{\star}), n(d^{\star}), L_0^{\star})$
\end{algorithmic}
\end{algorithm}

\section{Output characterization}

The returned design carries, beyond the geometry itself, a signed
error relative to the target,
\begin{equation}
\Delta n_f = n_f^{*}(d^{\star}) - n_f^{\star},
\label{eq:safety-error}
\end{equation}
with $\Delta n_f>0$ indicating a design more conservative than requested
and $\Delta n_f<0$ indicating a shortfall, together with the full
per-diameter candidate list (valid or not, with rejection reasons where
applicable) for traceability of the selection.

\section{Conclusion}

The inverse design problem — four unknown geometric parameters
constrained by one rate equation and one fatigue-safety target —
decomposes cleanly because the rate equation is linear in the coil count
(Eq.~\eqref{eq:active-coils}) and the fatigue constraint is monotonic in
the mean diameter for fixed wire diameter away from the $C=4$
singularity (Proposition~\ref{prop:monotonic}). This reduces a nominally
under-determined nonlinear system to a finite family of independent
scalar root-finding problems indexed by a discrete standard wire-diameter
series, each screened by simple geometric buildability inequalities
(Eqs.~\eqref{eq:min-coils}, \eqref{eq:no-bottom},
\eqref{eq:pitch-constraint}) before the closest-to-target candidate is
selected by Eq.~\eqref{eq:selection-rule}.

\end{document}
